% !TeX program = xelatex
\RequirePackage[T1]{fontenc}
\documentclass[journal,10pt]{IEEEtran}
\usepackage{amsmath,amssymb}
\usepackage{xeCJK}
\setmainfont{texgyretermes-regular.otf}[BoldFont=texgyretermes-bold.otf,ItalicFont=texgyretermes-italic.otf,BoldItalicFont=texgyretermes-bolditalic.otf]
\setCJKmainfont{FandolSong-Regular.otf}[BoldFont=FandolSong-Bold.otf,ItalicFont=FandolKai-Regular.otf,BoldItalicFont=FandolKai-Regular.otf,SmallCapsFont=FandolSong-Regular.otf]
\setCJKsansfont{FandolHei-Regular.otf}[BoldFont=FandolHei-Bold.otf]
\usepackage[unicode,hidelinks]{hyperref}
\newcommand{\E}{\mathbb E}
\newcommand{\Prob}{\mathbb P}
\newcommand{\R}{\mathbb R}
\newtheorem{proposition}{命题}
\renewcommand{\IEEEproofname}{证明}
\setlength{\emergencystretch}{1.5em}
\begin{document}
\title{面向拓扑辨识的需求响应设计}
\author{}
\maketitle
\begin{abstract}
给定当前模型和备选模型，需求响应应使两者的电压解释能力产生足以改变选模结果的差别。本文直接采用辨识器的损失改善条件，给出正态预测误差下满足该条件的概率保证，并将DR规划写为最小费用的二次约束问题。
\end{abstract}

\section{预测、损失与选模}
模型0表示当前选中的模型，模型1表示待比较的备选模型。两者均由历史数据拟合，之后固定参数，以$H_0,H_1\in\R^{m\times d}$将功率$P\in\R^d$映射为电压响应$V\in\R^m$。模型1可以增线、删线或改变连接，不预设两个模型之间的包含关系。

令$u\in\R^r$为DR计划调节量，功率预测为
\begin{equation}
P(u)\sim\mathcal N(\widehat p+Bu,\Sigma_P). \label{eq:power}
\end{equation}
其中$\widehat p$为无DR的功率预测均值，$B$将调节量转为节点P/Q变化，$\Sigma_P$为给定的功率预测误差协方差。电压误差采用共同模型
\begin{equation}
V\mid P,\text{模型}~k
\sim\mathcal N(H_kP,\Sigma_V),\quad k=0,1. \label{eq:voltage}
\end{equation}
$\Sigma_V\succ0$描述给定功率后剩余的电压误差，与功率预测误差独立。两项协方差均不随报价变化，可由预测器提供，也可作为假设给定。

用同一观测评价两个模型，定义平方损失
\begin{equation}
S_k=(V-H_kP)^\top\Sigma_V^{-1}(V-H_kP). \label{eq:loss}
\end{equation}
这里按电压误差大小对残差加权，并对全部观测求和。正态密度为
\begin{equation}
f_k(V\mid P)
=\frac{\exp(-S_k/2)}
{(2\pi)^{m/2}|\Sigma_V|^{1/2}}. \label{eq:density}
\end{equation}
设$q_h\in(0,1)$为历史更新后模型1的概率，$q_1$为加入新观测后的概率。两模型采用相同的功率预测分布，贝叶斯公式给出
\begin{equation}
\frac{q_1}{1-q_1}
=\frac{q_h}{1-q_h}\exp\!\left(\frac{S_0-S_1}{2}\right). \label{eq:bayes}
\end{equation}
式\eqref{eq:bayes}中的$1/2$来自正态密度\eqref{eq:density}。若把损失定义为$J_k=S_k/2$，指数即可写成$J_0-J_1$；相应的算法容差也应除以2。

辨识器采用的改善条件统一写为
\begin{equation}
S_0-\kappa S_1>\delta,\qquad \kappa\ge1. \label{eq:gate}
\end{equation}
$\kappa$由选模规则规定，$\delta$为新日需要补足的损失差。保留完整历史时，
\begin{equation}
\delta=\delta_{\rm all}-(S_0^h-\kappa S_1^h), \label{eq:history}
\end{equation}
其中$S_k^h$为历史损失，$\delta_{\rm all}$为合并数据的容差。平均损失应先按观测数换算到式\eqref{eq:loss}的总和口径。

规划时以模型1作为电压预测参考，即$V=H_1P+\varepsilon$、$\varepsilon\sim\mathcal N(0,\Sigma_V)$。以下$\Prob_1,\E_1$均指这一假设及给定历史下的概率和期望。要求
\begin{equation}
\Prob_1(S_0-\kappa S_1>\delta)\ge\beta,\qquad 0<\beta<1. \label{eq:chance}
\end{equation}
例如$\beta=0.9$表示在该预测模型下，至少90\%的执行满足改善条件。

\section{DR需要产生多大的电压差}
两个模型的预测电压之差及其协方差为
\begin{align}
d(u)&=(H_1-H_0)(\widehat p+Bu), \label{eq:d}\\
\Sigma_D&=(H_1-H_0)\Sigma_P(H_1-H_0)^\top. \label{eq:cov}
\end{align}
$d(u)$直接表示DR后两模型的平均预测差，$\Sigma_D$表示功率预测误差使该差异如何波动。

给定上述误差假设，预先计算固定权重
\begin{equation}
W=\frac{1}{4(\kappa+1)}
\left(\Sigma_V+\frac{\Sigma_D}{2(\kappa+1)}\right)^{-1} \label{eq:weight}
\end{equation}
及所需门槛
\begin{equation}
\begin{split}
h={}&\frac{\delta}{4\kappa}+\log\frac1{1-\beta}\\
&-\frac12\log
\frac{\left|\Sigma_V+\Sigma_D/[2(\kappa+1)]\right|}{|\Sigma_V|}\\
&-\frac m2\log\frac{\kappa+1}{2\kappa}.
\end{split} \label{eq:h}
\end{equation}
$W$按误差大小衡量电压差；$h$同时计入选模门槛和要求的成功比例。

\begin{proposition}
在式\eqref{eq:power}--\eqref{eq:chance}的固定参数模型下，
\begin{equation}
\boxed{d(u)^\top Wd(u)\ge h} \label{eq:condition}
\end{equation}
是式\eqref{eq:chance}的充分条件。
\end{proposition}
证明见附录。正文只需计算$d(u)$并与固定的$W,h$组合，无需在每次优化迭代中计算概率积分。

设$K(u)$为合同费用，$\mathcal U$为设备功率、能量和可用时段允许的调节范围。DR规划直接写为
\begin{equation}
\boxed{
\begin{aligned}
\min_{u\in\mathcal U}\quad &K(u),\\
\mathrm{s.t.}\quad &d(u)^\top Wd(u)\ge h.
\end{aligned}} \label{eq:optimization}
\end{equation}
$d(u)$对$u$为线性加常数，$W,h$在规划前固定，因此约束已经是二次形式，无需再引入一组展开系数。二次费用和线性设备约束对应非凸QCQP；原设备控制中的离散选择按原接口保留。

本节不新增矩估计步骤。均值和协方差是预测输入，证明中的期望在正态假设下直接计算。式\eqref{eq:condition}用一个可计算的上界控制失败概率，因此可能多要求一些调节量。它针对冻结模型的比较，重新调权与最终输出仍按实际辨识流程检验。

\clearpage
\appendices
\section{概率条件的推导}

\subsection{整理损失差}
记$\Delta H=H_1-H_0$。规划时$V=H_1P+\varepsilon$，
所以两模型的残差分别为
\[
V-H_1P=\varepsilon,\qquad
V-H_0P=\varepsilon+\Delta H P.
\]
为按同一电压误差尺度计算平方损失，令
\begin{equation}
e=\Sigma_V^{-1/2}\varepsilon,\qquad
Z=\Sigma_V^{-1/2}\Delta H P, \label{eq:app_ez}
\end{equation}
其中$\Sigma_V^{-1/2}$取对称正定平方根的逆。
由正文的正态模型，
\[
e\sim\mathcal N(0,I_m),\quad Z\sim\mathcal N(a,C),\quad e\perp Z,
\]
且
\[
a=\Sigma_V^{-1/2}d(u),\qquad
C=\Sigma_V^{-1/2}\Sigma_D\Sigma_V^{-1/2}.
\]
在给定$u$时，$a,C$均已知。于是
\begin{align}
S_1&=\|e\|^2,\qquad S_0=\|e+Z\|^2,\notag\\
L:=S_0-\kappa S_1
&=\|Z\|^2+2e^\top Z-(\kappa-1)\|e\|^2. \label{eq:app_L}
\end{align}
需要控制的失败事件是$L\le\delta$。

\subsection{选择概率上界}
直接计算$\Prob_1(L\le\delta)$，需要在式\eqref{eq:app_L}
给出的二次边界内积分；该边界还会随DR计划$u$改变。
为得到可预计算的二次约束，本文选择一个易于解析计算的概率上界。

由于$L$可能为负，先对任意$\lambda>0$定义非负随机量
$X=e^{-\lambda L}$。指数函数的单调性给出
\[
L\le\delta\quad\Longleftrightarrow\quad
X\ge e^{-\lambda\delta}.
\]
任意非负$X$与$\tau>0$满足
$\E_1X\ge\tau\Prob_1(X\ge\tau)$。
取$\tau=e^{-\lambda\delta}$，得到
\begin{equation}
\Prob_1(L\le\delta)
\le e^{\lambda\delta}\E_1e^{-\lambda L}. \label{eq:app_bound}
\end{equation}
因此，只需使右端不超过$1-\beta$。
使用式\eqref{eq:app_bound}是为保留可求解的二次形式所作的选择；
它给出充分条件，并非成功概率的精确表达。

\subsection{选择参数}
$\lambda$是证明和规划中由我们固定的数，不是待估计的负荷参数。
将正态密度中的$-\|e\|^2/2$与$-\lambda L$相加，
配方后会出现
\begin{equation}
b(\lambda)=1-2\lambda(\kappa-1),\qquad
\gamma(\lambda)
=\frac{\lambda(1-2\kappa\lambda)}{b(\lambda)}. \label{eq:app_general}
\end{equation}
取充分范围
\begin{equation}
0<\lambda<\frac1{2\kappa} \label{eq:app_range}
\end{equation}
可保证$b(\lambda)>0,\gamma(\lambda)>0$：
第一次积分有界，第二次积分对任意$C\succeq0$均有界，
并保留对模型电压差的正权重。式\eqref{eq:app_range}
是便于统一设计的充分范围，不是积分存在的必要范围。

本文采用固定值
\begin{equation}
\lambda=\frac1{4\kappa},\qquad
b=\frac{\kappa+1}{2\kappa},\qquad
\gamma=\frac1{4(\kappa+1)}. \label{eq:app_fixed}
\end{equation}
它位于式\eqref{eq:app_range}内，使后续公式简洁，
并避免在DR优化中再同时选择$\lambda$。
其他固定的合法取值也可形成二次约束；
式\eqref{eq:app_fixed}不保证上界最紧。

\subsection{计算右侧期望}
按全期望公式，可先固定$Z$、对电压误差$e$计算，
再对由功率误差引起的$Z$计算：
\begin{equation}
\E_1e^{-\lambda L}
=\E_Z\!\left[\E_e(e^{-\lambda L}\mid Z)\right]. \label{eq:app_tower}
\end{equation}
两次积分分别处理两种误差，不会重复计入同一误差。

\subsubsection{先对电压误差积分}
给定$Z=z$时，$z$是常数；被积变量只有$e$。
由于$e$与$Z$独立，$e$的条件密度仍为
\[
\phi_m(e)=(2\pi)^{-m/2}e^{-\|e\|^2/2}.
\]
期望就是函数乘以该密度后积分：
\begin{align}
\E_e(e^{-\lambda L}\mid Z=z)
&=\int_{\R^m}
 e^{-\lambda L(e,z)}\phi_m(e)\,de\notag\\
&=(2\pi)^{-m/2}e^{-\lambda\|z\|^2}\notag\\
&\quad\cdot\int_{\R^m}
 e^{-b\|e\|^2/2-2\lambda e^\top z}\,de. \label{eq:app_e_integral}
\end{align}
其中$b=1-2\lambda(\kappa-1)$正是合并正态密度后的平方项系数。
完成平方：
\begin{equation}
\begin{split}
-\frac b2\|e\|^2-2\lambda e^\top z
={}&-\frac b2\left\|e+\frac{2\lambda}{b}z\right\|^2\\
&+\frac{2\lambda^2}{b}\|z\|^2.
\end{split} \label{eq:app_e_square}
\end{equation}
令$y=\sqrt b\,(e+2\lambda z/b)$，
则$de=b^{-m/2}dy$，积分区域仍为$\R^m$。
利用标准正态密度的积分为1，即
\[
(2\pi)^{-m/2}\int_{\R^m}e^{-\|y\|^2/2}\,dy=1,
\]
式\eqref{eq:app_e_integral}便成为
\begin{equation}
\E_e(e^{-\lambda L}\mid Z=z)
=b^{-m/2}e^{-\gamma\|z\|^2}, \label{eq:app_e_result}
\end{equation}
这里$\gamma=\lambda-2\lambda^2/b$，
与式\eqref{eq:app_general}一致。
因子$b^{-m/2}$来自换元时的体积变化。

\subsubsection{再对功率误差积分}
由式\eqref{eq:app_tower}和式\eqref{eq:app_e_result}，
剩下的计算是
\begin{equation}
\E_1e^{-\lambda L}=b^{-m/2}\E_Z e^{-\gamma\|Z\|^2}. \label{eq:app_remaining}
\end{equation}
为同时允许$C$奇异，取$F\in\R^{m\times s}$使$FF^\top=C$，
写成$Z=a+F\xi$，其中$\xi\sim\mathcal N(0,I_s)$。
此时被积变量为$\xi$，$a,F$均固定。将函数与$\xi$的正态密度相乘：
\begin{align}
\E_Ze^{-\gamma\|Z\|^2}
={}&(2\pi)^{-s/2}\int_{\R^s}
e^{-\gamma\|a+F\xi\|^2-\|\xi\|^2/2}\,d\xi\notag\\
={}&(2\pi)^{-s/2}e^{-\gamma a^\top a}\notag\\
&\cdot\int_{\R^s}
e^{-\xi^\top D\xi/2-t^\top\xi}\,d\xi, \label{eq:app_xi_integral}
\end{align}
其中
\[
D=I_s+2\gamma F^\top F\succ0,\qquad t=2\gamma F^\top a.
\]
再次完成平方：
\begin{equation}
\begin{split}
-\tfrac12\xi^\top D\xi-t^\top\xi
={}&-\tfrac12(\xi+D^{-1}t)^\top D(\xi+D^{-1}t)\\
&+\tfrac12t^\top D^{-1}t.
\end{split}
\end{equation}
换元$w=D^{1/2}(\xi+D^{-1}t)$，有
$d\xi=|D|^{-1/2}dw$。
剩余标准正态密度的积分同样为1，故
\begin{equation}
\E_Ze^{-\gamma\|Z\|^2}
=|D|^{-1/2}
e^{-\gamma a^\top a+t^\top D^{-1}t/2}. \label{eq:app_xi_result}
\end{equation}
行列式因子$|D|^{-1/2}$同样来自换元，
不是额外加入的惩罚。

\subsection{整理为正文的二次条件}
令$M=I_m+2\gamma C$。利用
\[
|D|=|M|,\qquad
M^{-1}=I_m-2\gamma FD^{-1}F^\top,
\]
将式\eqref{eq:app_remaining}和式\eqref{eq:app_xi_result}合并，得到
\begin{equation}
\begin{split}
-\log\E_1e^{-\lambda L}
={}&\gamma a^\top M^{-1}a\\
&+\frac12\log|M|+\frac m2\log b.
\end{split} \label{eq:app_closed}
\end{equation}
上述两个矩阵恒等式分别由$|I+AB|=|I+BA|$
以及直接相乘验证。此处计算的只是给定正态分布下的期望，
没有新增参数估计。

令$G=\Sigma_V+2\gamma\Sigma_D$。由$a,C$的定义，
\[
M=\Sigma_V^{-1/2}G\Sigma_V^{-1/2}.
\]
因此
\[
\gamma a^\top M^{-1}a=d(u)^\top Wd(u),
\qquad |M|=\frac{|G|}{|\Sigma_V|},
\]
其中$W=\gamma G^{-1}$，正是正文式\eqref{eq:weight}。
式\eqref{eq:app_closed}可写为
\begin{equation}
-\log\E_1e^{-\lambda L}
=d(u)^\top Wd(u)
+\frac12\log\frac{|G|}{|\Sigma_V|}
+\frac m2\log b. \label{eq:app_physical}
\end{equation}
要求式\eqref{eq:app_bound}的右端不超过$1-\beta$，
取对数并代入式\eqref{eq:app_physical}，即
\begin{equation}
\begin{split}
d(u)^\top Wd(u)\ge{}&
\lambda\delta+\log\frac1{1-\beta}\\
&-\frac12\log\frac{|G|}{|\Sigma_V|}-\frac m2\log b.
\end{split} \label{eq:app_final}
\end{equation}
最后代入式\eqref{eq:app_fixed}，右端就是正文式\eqref{eq:h}的$h$，
从而得到$d(u)^\top Wd(u)\ge h$。
正态积分的计算是精确的；概率保证的保守性来自所选上界及$\lambda$的取值。
\end{document}
